% Output-guided behavioral decomposition: the method as implemented in ace/.
%
% Drop-in replacement for the earlier section. Every step below describes what the code in
% `ace/` actually does; file and line references are given in comments so the two can be
% checked against each other. Requires the macros \ACE, \Tr, \Grammar, \Regions and the
% packages already used by the paper (amsmath, algorithm, algorithmicx/algpseudocode, tikz).

\subsection{Output-event labeling}
\label{sec:labeling}
For each selected event predicate $e$, \ACE{} collects the occurrence set
\begin{equation}
E_e(\Tr)=\{(k,t)\mid \tau_k,t\models_f e\},
\end{equation}
by a single linear scan of the original traces.\footnote{\texttt{ace/labeling.py:41--54}.}
This replaces the former output stuck-at wrapper: for a Boolean output, forcing the opposite
constant and comparing against the original trace yields a label determined entirely by the
already observed output value, so the wrapper, the re-simulation it required and the
synthetic detection signal it produced are all redundant. Direct labeling preserves the
useful anchoring idea without claiming fault propagation or internal cone discovery. The IP
is never instrumented, encapsulated, forced or re-simulated; the cost of the step is
$O(|\Tr|)$ predicate evaluations and no simulator time at all.

Two conventions matter for what follows.

\emph{Occurrences are recorded where the match ends.} A multi-sample event has an extent;
$\ACE$ indexes it by its last sample, because trigger scoring (\S\ref{sec:mining}) reads the
occurrences it explains off the same match-end positions. Indexing by the first sample
instead makes the two sets disjoint, every candidate appears to explain nothing, and the
region comes out empty.

\emph{A held level is a poor event, and is rewritten as its onset.} An output flag that
stays high between requests is true at a large fraction of samples; then every position has
the event within the horizon, no trigger can beat that base rate, and the entire response
family is discarded for want of a candidate. Setting \texttt{event\_onsets} rewrites each
declared event as
\begin{equation}
e^{\uparrow}\equiv(\lnot e)\,\texttt{\#\#1}\,e ,
\end{equation}
which turns the 2189 samples on which \texttt{done} is high in the square-root corpus into
247 completions.\footnote{\texttt{ace/\_\_main\_\_.py:54--65}.} Signals that are already
pulses are unaffected, since for them every high sample is a rising edge.

In the running example, choose $e_{done}(t)\equiv[\texttt{done}(t)=1]$ under the onset
rewriting: the event positions are the cycles in which \texttt{done} rises. A second event
$e_{err}(t)\equiv[\texttt{error}(t)=1]$ anchors the error behavior. Nothing restricts events
to single output bits: for multi-bit outputs and protocol interfaces the event vocabulary
supplies transitions, thresholds or handshake milestones (\texttt{pready == 1},
\texttt{rd\_data\_o >= 8388608}), so the same machinery applies without equating each output
value with a functionality.

\subsection{Mining explanatory candidates}
\label{sec:candidates}
Given $e$, the temporal-property backend enumerates candidates $a\in\Grammar$ that can
predict $e$ under temporal horizon $H$. \ACE{} obtains them from HARM by pinning the event
into the last proposition slot of every template of the configured grammar, mining, and
taking the deduplicated antecedents of what comes back.\footnote{\texttt{ace/\_\_main\_\_.py:84--113};
grammars \texttt{G1}--\texttt{G5} in \texttt{ace/backends.py:56--72}, default \texttt{G3}
$=\{G(P_0\!\mid\!\to P_1),\,G(P_0\!\mid\!\Rightarrow P_1),\,G(P_0\!\mid\!\to\,\texttt{\#\#[1:H]}\,P_1)\}$.}
The miner is therefore used as a candidate generator, not as an oracle: metrics are never
read back from it, since \ACE{} scores every candidate with its own finite-trace evaluator,
which is what keeps labelling, trigger selection and held-out validation in one semantics.
The mined antecedents are unioned with a vocabulary \ACE{} reads off the traces itself --
the observed values, the interval bounds, the three orderings $a>b$, $a=b$, $a<b$ for every
input pair and their pairwise conjunctions.\footnote{\texttt{ace/\_\_main\_\_.py:116--140}.}
The union is not redundant: value clustering only ever yields comparisons against constants,
so a relation between two signals -- what explains a comparator asserting its output -- is
outside the miner's vocabulary by construction.

\subsection{Scoring candidates}
\label{sec:mining}
For a candidate $a$, the candidate assertion is $G(a\mid\to\,\texttt{\#\#[1:H]}\,e)$, and
\ACE{} scores it over the \emph{event occurrences}:
\begin{align}
ATCT(a,e) &= |E_{a,e}|,  &&\text{occurrences of } e \text{ that } a \text{ explains},\\
AFCT(a,e) &= |E_e(\Tr)| - |E_{a,e}|, &&\text{occurrences it leaves unexplained},
\end{align}
so $ATCT$ is $TP$, $AFCT$ is $FN$, and $ATCT+AFCT=|E_e(\Tr)|$ for every candidate of one
event. An occurrence is explained when $a$ matches at a position from which that occurrence
falls inside the horizon. A match counts only where $a$ was false at the previous sample:
without this, a request held high for a whole operation contributes one match per held cycle,
and on a 36-cycle divider whose enable is held throughout, every candidate scores the base
rate and the design yields no region at all. Write $m(a)$ for that onset-guarded match count,
which is the candidate's statistical support and is \emph{not} $ATCT+AFCT$.

Candidate quality is summarized by the smoothed recall
\begin{equation}
\hat{r}(a)=\frac{ATCT(a,e)+1}{ATCT(a,e)+AFCT(a,e)+2},
\label{eq:recall}
\end{equation}
where the $+1/+2$ smoothing penalises candidates with little statistical support and keeps
the number comparable across traces of different length, which a raw $ATCT-AFCT$ volume score
cannot do. \ACE{} also records, per candidate, the precision $|{\rm answered\ firings}|/m(a)$
and the full per-sample contingency table in the miner's own naming; neither is used to
select, both are emitted so any other figure of merit can be recomputed from the artifact.

\emph{A note on naming.} A temporal miner's contingency table is indexed by \emph{samples},
where \texttt{afct} denotes antecedent-false and consequent-true --- the event occurred and
the antecedent did not predict it. That is the same cell as $AFCT$ above, which is why the
ranking expression handed to the backend is literally
Eq.~\ref{eq:recall}.\footnote{\texttt{ace/backends.py}, \texttt{ace/triggers.py:47--79}.} The
sample-indexed table also has an \texttt{atcf} cell --- the antecedent fired and the event did
not follow --- which is the false positive and has no counterpart in an occurrence-indexed
population, since every occurrence of $e$ has the consequent true by construction.

Recall alone would rank a predicate that fires constantly first: matching at every second
sample is followed by every occurrence of $e$ without predicting any of them. \ACE{} therefore
compares each candidate against a null model built on its own match count. With $b$ the share
of positions from which $e$ follows within $H$,
\begin{equation}
\hat{r}_0(m)=\frac{\min(b\,m,\ |E_e(\Tr)|)+1}{|E_e(\Tr)|+2},\qquad
\mathit{lift}(a)=\hat{r}(a)-\hat{r}_0(m(a)),
\label{eq:lift}
\end{equation}
i.e.\ what a predicate firing $m$ times at arbitrary positions would be expected to explain,
capped because a firing cannot explain an occurrence twice and smoothed the same way as the
candidate. A candidate that fires thousands of times saturates the expectation and its lift
falls to zero; a trigger that fires once per operation has a small expectation and a large
lift.

The test is deliberately strict, and it is strictest when the declared horizon exceeds the
real response latency: if every position in the corpus has an occurrence of $e$ within $H$,
no predicate can beat chance and the region is reported unexplained rather than anchored on a
coincidence. The event coverage of the \emph{selected set},
\begin{equation}
\mathit{cov}(P_e)=\frac{|E_{P_e,e}|}{|E_e(\Tr)|},
\label{eq:cov}
\end{equation}
is reported alongside and is the first term of the selection objective below.

\subsection{Selecting a trigger set}
Algorithm~\ref{alg:decomp} makes the selection policy explicit; it is implemented in
\texttt{ace/triggers.py:221--307} and the thresholds below are its configuration keys. The
filter phase rejects a candidate that never matched, that matched too rarely to explain the
event ($m(a)$ below a fraction $\sigma$ of $|E|$, which separates ``explains the event'' from
``was true when it happened''), that fails to beat its null model by $\theta_{lift}$, or that
falls below a plain recall floor $\theta_{r}$. That floor is deliberately low: it asks what a
\emph{single} trigger explains, and a region with two complementary triggers has no candidate
above one half by construction. Every rejection is recorded with its reason and reaches the
output artifact, so nothing is dropped silently.

The greedy phase maximises
\begin{equation}
score(a)=\mathit{gain}(a)+\mathit{lift}(a)-\lambda\,\mathit{red}(a),
\label{eq:score}
\end{equation}
with $\mathit{gain}(a)=|E_{a,e}\cap U|/|E|$ the marginal coverage of the still-unexplained
occurrences $U$ and $\mathit{red}(a)$ the largest Jaccard overlap between the occurrences $a$
explains and those already explained by a selected trigger. Two points of the implementation
are worth stating plainly. First, $\theta_{gain}$ thresholds the \emph{combined} score rather
than the gain term alone, so a candidate with modest marginal coverage but high lift can
still be admitted; the loop terminates when no candidate clears the threshold, when every
occurrence is explained, or at $k_{max}$ triggers. Second, the second term is the lift of
Eq.~\ref{eq:lift} and not the raw recall of Eq.~\ref{eq:recall}, because a predicate that
fires constantly has recall $1$ and explains nothing. Marginal coverage drives the loop, which is what
brings complementary triggers out together — a multiplier enable and a divider enable, each
accounting for its own half of \texttt{valid\_o}.

\begin{algorithm}[t]
\caption{Output-Guided Behavioral Decomposition}
\label{alg:decomp}
\scriptsize
\begin{algorithmic}[1]
\Require Trace set $\Tr$; event predicate $e$; grammar $\Grammar$; horizon $H$; thresholds $\theta_{r},\theta_{lift},\theta_{gain}$; support fraction $\sigma$; redundancy weight $\lambda$; trigger budget $k_{max}$; window sizes $h_{pre},h_{post}$
\Ensure Trigger set $P_e$; behavioral regions $\Regions_e$; episode sets $\{\Tr_r\}$
\State $E\gets E_e(\Tr)$; $U\gets E$; $P_e\gets\emptyset$ \Comment{\S\ref{sec:labeling}}
\State $C\gets\textsc{MineCandidates}(\Tr,e,\Grammar,H)$ \Comment{\S\ref{sec:candidates}}
\State $b\gets\textsc{BaseRate}(\Tr,e,H)$
\ForAll{$a\in C$}
  \State compute $E_{a,e}$, $m(a)$, $ATCT(a,e)$, $AFCT(a,e)$, $\hat{r}(a)$, $\mathit{lift}(a)$
\EndFor
\State $C\gets\{a\in C\mid m(a)\ge\max(2,\sigma|E|)\ \wedge\ \mathit{lift}(a)\ge\theta_{lift}\ \wedge\ \hat{r}(a)\ge\theta_{r}\}$ \label{alg:filter}
\While{$U\neq\emptyset$ \textbf{and} $|P_e|<k_{max}$} \label{alg:greedy-start}
  \ForAll{$a\in C\setminus P_e$}
    \State $gain(a)\gets |E_{a,e}\cap U|/|E|$
    \State $red(a)\gets\max_{p\in P_e}\textsc{Jaccard}(E_{a,e},E_{p,e})$
    \State $score(a)\gets gain(a)+\mathit{lift}(a)-\lambda\,red(a)$
  \EndFor
  \State $a^*\gets\arg\max score(a)$
  \If{$score(a^*)\le\theta_{gain}$} \State \textbf{break} \EndIf
  \State $P_e\gets P_e\cup\{a^*\}$; $U\gets U\setminus E_{a^*,e}$
\EndWhile \label{alg:greedy-end}
\State $A\gets\{(k,t)\in E\mid (k,t)\text{ explained}\}\cup\bigcup_{a\in P_e}\text{answered firings}(a)$ \label{alg:anchors}
\State $\Tr_r\gets\textsc{ExtractEpisodes}(\Tr,A,h_{pre},h_{post})$
\State $r\gets(e,P_e,\Tr_r)$; $\Regions_e\gets\Regions_e\cup\{r\}$
\State \Return $P_e,\Regions_e,\{\Tr_r\}$, $U$ \Comment{$U$: occurrences left unassigned}
\end{algorithmic}
\end{algorithm}

The selection is coverage-oriented rather than requiring complete and unique coverage. Two
triggers explaining the same occurrences are retained only if they are not semantically
redundant, which is what the $\lambda$ term measures. Disjoint triggers naturally create
distinct regions. If no candidate passes the filter of Line~\ref{alg:filter} or the
marginal-score threshold, the corresponding event occurrences remain explicitly unassigned
rather than being forced into a spurious region, and the region as a whole may be reported
as skipped.

In the square-root example, a high-quality trigger for $e_{done}$ captures a two-cycle
relation from a valid request, e.g.\ a predicate equivalent to ``$\texttt{start}=1$ and
$\texttt{in}\ge0$ two cycles before $\texttt{done}=1$''. This trigger induces the
normal-completion region. A separate trigger for $e_{err}$ captures the negative-input error
region.

\subsection{Boundary-preserving episode extraction}
The region of a trigger set is anchored at two kinds of position (Line~\ref{alg:anchors}):
the event occurrences that the selected triggers explain, and the $ATCT$ positions of those
triggers. Both are needed, because an episode must contain the stimulus \emph{and} the
response; matches with no response are deliberately excluded, being evidence against the
trigger rather than a region to mine.

For each anchor $(k,t)$, \ACE{} extracts the contiguous window $[t-h_{pre},t+h_{post}]$ from
the same $\tau_k$, clipped at that run's boundaries.\footnote{\texttt{ace/episodes.py:37--50}.}
All original input and output values are preserved. Windows that overlap or touch are merged,
and since windows are accumulated per run index, merging across two executions is not merely
forbidden but unrepresentable. The window sizes are not padding: a temporal miner cannot
deduce cause and effect from one isolated sample, so $h_{pre}$ exposes the setup conditions
before the trigger fires and $h_{post}$ the obligations after the event. If $h_{post}<H$ the
flow warns and records that bounded-response clauses cannot be decided inside an episode,
since the response window would leave the episode and every such guarantee would be
unreachable by construction.

\textbf{Proposition 1 (episode provenance).} Every episode produced by the extraction step is
a contiguous subsequence of exactly one original execution in $\Tr$.

\emph{Proof sketch.} Episode construction selects a bounded index interval from one $\tau_k$,
clips it at the trace endpoints, and performs neither value substitution nor concatenation
across traces. \hfill$\square$

The proposition is not only argued but checked at run time: the predicate
$0\le lo\le hi<|\tau_k|$ is verified for every episode and its result is written into the
region manifest and the output artifact, so a violation would be visible rather than
silent.\footnote{\texttt{ace/episodes.py:81--83}.} Each episode is then re-exported as a run
of its own, which is what makes an episode boundary semantically real: a delay reaching past
the end of an episode simply fails to match instead of silently reading the neighbouring
samples of the original execution. Episodes are written one CSV per episode, so that no
backend can relate samples across a boundary. A second mode concatenates a region's episodes
into a single file for speed; it admits matches across an episode seam and is therefore
reserved for horizons small relative to $h_{pre}+h_{post}$. All results reported here use the
per-episode mode, which is the default.

\subsection{Cross-event merging}
Events are declared, and two declarations can name one behavior: a level and its negation, or
a handshake milestone and the transfer it completes. Since each event is decomposed
independently, \ACE{} closes the flow with one pass over the finished regions and groups those
that are the same region. Two regions are grouped when their trigger sets are equivalent on
the observed traces \emph{and} their guarantee sets are equivalent, where equivalence between
two sets of clauses means that every member of one is entailed, over the corpus, by some
member of the other and conversely.\footnote{\texttt{ace/\_\_main\_\_.py:149--193}, reusing the
trace-bounded entailment check of \texttt{ace/validation.py:80}.} Requiring both halves is
what keeps two events that happen to share a trigger but impose different obligations in
separate regions — the paper's criterion that redundant triggers are merged only when their
downstream regions do not produce meaningfully different contracts. The pass is additive:
groups are reported, and the per-event regions, their episodes and their contracts are left
intact for inspection.

\begin{figure}[t]
\centering
\begin{tikzpicture}[x=0.43cm,y=0.38cm,font=\scriptsize,>=Latex]
% axes
\draw[->] (0,0)--(13,0) node[right]{time};
\node[left] at (0,2.9) {\texttt{start}};
\node[left] at (0,2.0) {\texttt{in}$\ge0$};
\node[left] at (0,1.1) {\texttt{done}};
\draw (0,2.7)--(1,2.7)--(1,3.2)--(2,3.2)--(2,2.7)--(5,2.7)--(5,3.2)--(6,3.2)--(6,2.7)--(13,2.7);
\draw (0,1.8)--(4,1.8)--(4,2.3)--(8,2.3)--(8,1.8)--(13,1.8);
\draw (0,0.9)--(3,0.9)--(3,1.4)--(4,1.4)--(4,0.9)--(7,0.9)--(7,1.4)--(8,1.4)--(8,0.9)--(13,0.9);
\draw[dashed] (3,0.4)--(3,3.5); \draw[dashed] (7,0.4)--(7,3.5);
\node[draw,rounded corners,fill=gray!5,align=left] at (6.4,-1.8) {event: $e_{done}\equiv$ rising edge of $done$\\trigger: valid request $\rightarrow$ $e_{done}$ after 2 cycles, counted at its onset\\anchors: the explained events and the request onsets that were answered\\region: original episodes $[t-h_{pre},t+h_{post}]$ around those anchors\\candidate $A$: $in\ge0$ (plus temporal environment clauses, if mined)\\candidate $G$: request $\rightarrow$ completion after 2 cycles};
\draw[->] (3,0.5)--(4.2,-0.6); \draw[->] (7,0.5)--(7.7,-0.6);
\end{tikzpicture}
\caption{Running example of the proposed flow. Output events anchor trigger mining; the
resulting region contains only original, boundary-preserving trace episodes. The final
contract is validated semantically against the documented square-root behavior.}
\label{fig:example}
\end{figure}

\subsection{What decomposition does and does not guarantee}
A region states that, in the available observations, a trigger explains an output event with
measured recall and measured precision. It does not prove causal uniqueness, expose an
internal logic cone, or guarantee that all implementation behaviors have been observed. The
recall of Eq.~\ref{eq:recall} is an association on the observed corpus, compared against a
null model at the candidate's own match count; it is not a measure of causal strength, and a
trigger can reach a high recall by firing often enough that the event follows it by
coincidence --- which is precisely what Eq.~\ref{eq:lift} exists to detect and what leaves
some events unexplained rather than explained badly. Region overlap, unassigned occurrences,
rejected candidates with their reasons, the per-sample contingency table of every candidate,
and episode provenance are all measured and emitted rather than hidden.
